\input{../tex-inputs/latex-leseansicht-vorspann}

\section[Theodor von Sosnosky an Arthur Schnitzler, 3. 2. 1902]{L04322 Theodor von Sosnosky an Arthur Schnitzler, 3. 2. 1902}
\nopagebreak\mylabel{L04322v}
\rehead{ }\normalsize\beginnumbering\briefempfaengerindex{Schnitzler, Arthur@\textsc{Schnitzler, Arthur}!zzzSosnosky, Theodor von@\emph{von Theodor von Sosnosky}!1902-03-031@{3. 2. 1902}|(be}
\toendnotes[C]{\smallbreak\pagebreak[2]}
\correspDesc{Versand  durch Theodor von Sosnosky am 3. 2. 1902 in Alland
\newline{}Erhalt  durch Arthur Schnitzler im Zeitraum [4. 3. 1902
                  – 8. 3. 1902?] in Wien}\toendnotes[C]{\smallbreak}
\Standort{DLA, A:Schnitzler, HS.1985.1.4640.}
\physDesc{Brief, 1 Blatt, 2 Seiten, 840 Zeichen
\newline{}Handschrift: blaue Tinte, deutsche Kurrent
\newline{}Schnitzler: 1) beschriftet: »Sosnoski«  2) mit rotem Buntstift eine Unterstreichung}\toendnotes[C]{\smallbreak}
\pstart
           \raggedleft{}{\pb}Alland b Baden\oindex{Alland@\textbf{Alland}, \emph{Verwaltungsgebiet}|pw}\pend
           
\pstart
           \raggedleft{}3./III. 902\pend
           
\pstart{}Sehr geehrter Herr Doktor!\pend\vspace{0.5em}
\pstart
           Anbei erlaube ich mir, Ihnen \label{K_L04322-1v}\edtext{unter X Band}{\lemma{\textnormal{\emph{unter X Band}}}\Cendnote{\textnormal{Unter Kreuzband, zwei sich rechtwinklig kreuzenden Papier- oder
                  Pappestreifen, ließen sich Drucksachen zu ermäßigtem Tarif versenden.}}}\label{K_L04322-1}{ }\label{K_L04322-2v}\edtext{eine Beſprechung\pwindex{Sosnosky, Theodor von 4.\,1.\,1866 Budapest – 3.\,2.\,1943 Wien@\textsc{Sosnosky, Theodor von} (4.\,1.\,1866 Budapest – 3.\,2.\,1943 Wien), \emph{Schriftsteller}!Bücher und Zeitschriftenschau [Frau Bertha Garlan]@\strich\emph{Bücher und Zeitschriftenschau [Frau Bertha Garlan]}|pwv}}{\lemma{\textnormal{\emph{eine Besprechung}}}\Cendnote{\textnormal{Th. v. S.\pwindex{Sosnosky, Theodor von 4.\,1.\,1866 Budapest – 3.\,2.\,1943 Wien@\textsc{Sosnosky, Theodor von} (4.\,1.\,1866 Budapest – 3.\,2.\,1943 Wien), \emph{Schriftsteller}|pwk}: \emph{Bücher und Zeitschriftenschau}\pwindex{Sosnosky, Theodor von 4.\,1.\,1866 Budapest – 3.\,2.\,1943 Wien@\textsc{Sosnosky, Theodor von} (4.\,1.\,1866 Budapest – 3.\,2.\,1943 Wien), \emph{Schriftsteller}!Bücher und Zeitschriftenschau [Frau Bertha Garlan]@\strich\emph{Bücher und Zeitschriftenschau [Frau Bertha Garlan]}|pwk}. In: \emph{Norddeutsche Allgemeine Zeitung}\pwindex{Norddeutsche Allgemeine Zeitung@\emph{Norddeutsche Allgemeine Zeitung}|pwk}, Nr. 50,
                        28. 2. 1902, Beilage.}}}\label{K_L04322-2} Ihres Romans »Frau Bertha Garlan\pwindex{Schnitzler, Arthur 15. 5. 1862 Wien – 21. 10. 1931 ebd.@\textsc{Schnitzler, Arthur} (15. 5. 1862 Wien – 21. 10. 1931 ebd.), \emph{Schriftsteller, Mediziner}!Frau Bertha Garlan. Roman@\strich\emph{Frau Bertha Garlan. Roman}|pw}« zu überſenden. Da ſie ſo
               ungeheuerlich \label{K_L04322-3v}\edtext{\textsc{post festum}}{\lemma{\textnormal{\emph{post festum}}}\Cendnote{\textnormal{lateinisch: nach dem Fest, im
                  Nachhinein}}}\label{K_L04322-3} ko{\geminationm}t, ſo kann ich ſie nicht einfach
               wortlos an Sie{ }ſenden, ſondern fühle mich verpflichtet, dieſe ſonderbare Sendung zu
                  ko{\geminationm}entieren. Als Sie im Vorjahre die
               Freundlichkeit hatten, mir das Buch\pwindex{Schnitzler, Arthur 15. 5. 1862 Wien – 21. 10. 1931 ebd.@\textsc{Schnitzler, Arthur} (15. 5. 1862 Wien – 21. 10. 1931 ebd.), \emph{Schriftsteller, Mediziner}!Frau Bertha Garlan. Roman@\strich\emph{Frau Bertha Garlan. Roman}|pwv} zu ſenden, beſprach  ich es zweimal: \label{K_L04322-4v}\edtext{die eine Rezension\pwindex{Sosnosky, Theodor von 4.\,1.\,1866 Budapest – 3.\,2.\,1943 Wien@\textsc{Sosnosky, Theodor von} (4.\,1.\,1866 Budapest – 3.\,2.\,1943 Wien), \emph{Schriftsteller}!Neue Erzählungsliteratur@\strich\emph{Neue Erzählungsliteratur}|pwv}}{\lemma{\textnormal{\emph{die eine Rezension}}}\Cendnote{\textnormal{Theodor von Sosnosky\pwindex{Sosnosky, Theodor von 4.\,1.\,1866 Budapest – 3.\,2.\,1943 Wien@\textsc{Sosnosky, Theodor von} (4.\,1.\,1866 Budapest – 3.\,2.\,1943 Wien), \emph{Schriftsteller}|pwk}: \emph{Neue Erzählungsliteratur}\pwindex{Sosnosky, Theodor von 4.\,1.\,1866 Budapest – 3.\,2.\,1943 Wien@\textsc{Sosnosky, Theodor von} (4.\,1.\,1866 Budapest – 3.\,2.\,1943 Wien), \emph{Schriftsteller}!Neue Erzählungsliteratur@\strich\emph{Neue Erzählungsliteratur}|pwk}. In: \emph{Neues Wiener Tagblatt}\pwindex{Neues Wiener Tagblatt@\emph{Neues Wiener Tagblatt}|pwk}, Jg. 35, Nr.  167,
                     20. 1. 1901, S. [1]–3. Es handelt sich um eine
                  Sammelrezension, die außer \emph{Bertha Garlan}\pwindex{Schnitzler, Arthur 15. 5. 1862 Wien – 21. 10. 1931 ebd.@\textsc{Schnitzler, Arthur} (15. 5. 1862 Wien – 21. 10. 1931 ebd.), \emph{Schriftsteller, Mediziner}!Frau Bertha Garlan. Roman@\strich\emph{Frau Bertha Garlan. Roman}|pwk}
                  verschiedene Neuerscheinungen von Georg von
                     Ompteda\pwindex{Ompteda, Georg von 29.\,3.\,1863 Hannover – 10.\,12.\,1931 München@\textsc{Ompteda, Georg von} (29.\,3.\,1863 Hannover – 10.\,12.\,1931 München), \emph{Schriftsteller}|pwk}, Maximilian von Rosenberg\pwindex{Rosenberg, Maximilian von 1849 Halberstadt – 1913@\textsc{Rosenberg, Maximilian von} (1849 Halberstadt – 1913)|pwk},
                  Wolf von Tainach\pwindex{Tainach, Wolf von @\textsc{Tainach, Wolf von}, \emph{Schriftsteller}|pwk}, Emil Marriot\pwindex{Marriot, Emil 20.\,11.\,1855 Wien – 1938 ebd.@\textsc{Marriot, Emil} (20.\,11.\,1855 Wien – 1938 ebd.), \emph{Schriftstellerin}|pwk} und Paul von
                     Schönthan\pwindex{Schönthan-Pernwald, Paul von 19.\,3.\,1853 Wien – 4.\,8.\,1905 ebd.@\textsc{Schönthan-Pernwald, Paul von} (19.\,3.\,1853 Wien – 4.\,8.\,1905 ebd.), \emph{Schriftsteller, Journalist}|pwk} behandelt.}}}\label{K_L04322-4} erſchien im
                  Juni{ }\introOben{}1901\introOben{} mit andren Be{\pb}ſprechungen im Feuilleton
               des N. Wnr. Tagebl.\pwindex{Neues Wiener Tagblatt@\emph{Neues Wiener Tagblatt}|pw}\footnote{\noindent{}ich habe ſie Ihnen ſeinerzeit direkt \label{K_L04322-5v}\toendnotes[C]{\begin{minipage}[t]{4em}{\makebox[3.6em][r]{\tiny{Fußnote}}}\end{minipage}\begin{minipage}[t]{\dimexpr\linewidth-4em}\textit{zugeschickt}\,{]} Siehe XXXX Auszeichnungsfehler: Dokument L04232 nicht gefunden.\end{minipage}\par}zugeſchickt\label{1}.}, die andere\pwindex{Sosnosky, Theodor von 4.\,1.\,1866 Budapest – 3.\,2.\,1943 Wien@\textsc{Sosnosky, Theodor von} (4.\,1.\,1866 Budapest – 3.\,2.\,1943 Wien), \emph{Schriftsteller}!Bücher und Zeitschriftenschau [Frau Bertha Garlan]@\strich\emph{Bücher und Zeitschriftenschau [Frau Bertha Garlan]}|pwv} ſandte
               ich an die »Norddeutſche Allgemeine\pwindex{Norddeutsche Allgemeine Zeitung@\emph{Norddeutsche Allgemeine Zeitung}|pw}«, die{ }ſie
               zu meiner \uline{eigenen} Überraſchung \uline{heute} (!) zuſandte.\pend
           
\pstart
           Ich bitte also nicht zu glauben, daß \uline{ich} ſo ſaumſelig geweſen bin: \strikeout{mi} dies zu
               konſtatiren iſt der Zweck dieſes Briefes.\pend
           
\pstart
           Mit der verbindlichſten Empfehlung zeichne ich, geehrter Herr Doktor,{\\[\baselineskip]}Hochachtungsvoll{\\[\baselineskip]}\spacefill\mbox{Theodor vSosnosky}\pend
           \leftskip=0em{}\selectlanguage{ngerman}\endnumbering\briefempfaengerindex{Schnitzler, Arthur@\textsc{Schnitzler, Arthur}!zzzSosnosky, Theodor von@\emph{von Theodor von Sosnosky}!1902-03-031@{3. 2. 1902}|)be}\mylabel{L04322h}
\begin{anhang}
\end{anhang}\newcommand{\dateiname}{L04322}\newcommand{\titel}{Theodor von Sosnosky an Arthur Schnitzler, 3. 2. 1902}\newcommand{\editorInnen}{Herausgegeben von Jahnke, SelmaMüller, Martin Anton}\input{../tex-inputs/latex-leseansicht-abspann}
